\documentclass[11pt,a4paper]{article}
\usepackage{curso}
\begin{document}
\portada{Ejercicios --- Lección}{8.7}{Entrenamiento consciente del hardware}
\section{Ejercicios}
\begin{enumerate}[label=\textbf{E\arabic*.}]
\item (B) ¿Qué hace el estimador directo en el paso hacia atrás?
\item (B) ¿Por qué la inyección de ruido favorece soluciones planas?
\item (B) ¿Por qué fijar la máscara de dispersión desde el paso 0?
\item (B) Diseñe la evaluación de un método consciente del hardware.
\item (I) Para un peso $w$ con ruido $w(1+\sigma\xi)$, calcule el aumento esperado de una pérdida cuadrática $\frac12(wx-y)^2$ y relaciónelo con un término de regularización.
\item (I) ¿Cuántos niveles tiene un cuantizador simétrico de 4 bits sin el valor $-8$? ¿Cuál es el error máximo de redondeo relativo al peso máximo?
\item (A) Repita el experimento con 2, 3 y 8 bits y con $\sigma\in\{0,\,0{,}1,\,0{,}3\}$; grafique la ganancia del entrenamiento consciente.
\end{enumerate}
\section{Soluciones}
\begin{solucion}
\textbf{E1.} Trata el redondeo como identidad: pasa el gradiente respecto de $\hat W$ a la copia en coma flotante $W$.

\textbf{E2.} Minimiza la pérdida esperada bajo perturbaciones: penaliza la curvatura (soluciones donde pequeños cambios de los pesos no cambian la salida).

\textbf{E3.} Los pesos se adaptan a la conectividad real; podar después elimina pesos que la red usaba sin darle ocasión de compensar.

\textbf{E4.} Precisión en hardware simulado (bits, ruido, máscara) frente al despliegue ingenuo, varias semillas y varios sorteos de ruido.

\textbf{E5.} $\mathbb E[\frac12(w(1+\sigma\xi)x-y)^2]=\frac12(wx-y)^2+\frac12\sigma^2w^2x^2$: regularización $L_2$ ponderada por $x^2$.

\textbf{E6.} 15 niveles ($-7\dots7$); error máximo $\frac12\cdot\frac17\approx7\,\%$ del peso máximo.

\textbf{E7.} La ganancia crece cuando el hardware es más agresivo (pocos bits, mucho ruido).
\end{solucion}
\end{document}
